\scanpage{4}
% archive foliation: 3
Relation entre morphismes de topos et recollements de topos.

\margin{On remplace un morphisme \unsure{quelc.} de topos par un plongement fermé.}

1) \deleted{Soit}\inserted{Soit}
\[
(1)\qquad f:\mathcal X_0\longrightarrow\mathcal X_1\quad\deleted{\unsure{\mathcal X}}
\]
un morphisme de topos, on peut le considérer comme un topos fibré sur $\Delta_1$ (fibres $\mathcal X_0,\mathcal X_1$), et regarder le topos total associé $\mathcal X=C(f)$ \inserted{(le mapping-cone)} dont les faisceaux sont les systèmes
\[
(2)\qquad\left\{\begin{array}{ll}
(F_0,F_1,u),&F_0\in\operatorname{Ob}\mathcal F(\mathcal X_0),\ F_1\in\operatorname{Ob}\mathcal F(\mathcal X_1)\\
&u:\unsure{f_0^*|f^*}(F_1)\longrightarrow F_0\quad\text{i.e.}\\
&F_1\longrightarrow f_*(F_0)
\end{array}\right.
\]
\margin{\uline{NB} On a $\mathcal X\longrightarrow\operatorname{Top}(\Delta_1)$, ou ce qui revient au $=$, on a un $\mathcal X_0$ \unsure{comme} sous-topos ouvert (image inverse du sous-objet $\{0\}$ de l'objet final de $\widehat{\Delta}_1$).}

Alors $\mathcal X_0$ peut être regardé comme un ouvert de $\mathcal X$, $\mathcal X_1$ comme le fermé complémentaire, mais de plus l'inclusion
\[
(3)\qquad\mathcal X_1\xrightarrow{i_1}\mathcal X
\]
\uline{est strictement \deleted{asphérique}\inserted{\unsure{équiv. f.|équiv. h.}}}
i.e. pour \uline{tout} faisceau $F$ \inserted{$(F_0,F_1,u)$} sur $\mathcal X$,
\[
H^*(\mathcal X,F)\xrightarrow{\sim}H^*(\mathcal X_1,i_1^*(F)).
\]
\margin{\uline{Re NB} On a une rétraction $r$ de $\mathcal X$ sur le sous-topos fermé $\mathcal X_1$, par le foncteur
\[
r^*(F_1)=(f^*F_1\unsure{)},F_1,\deleted{\unsure{Ident.}}).
\]
\[
\mathcal X=(\mathcal X_0\times\widehat{\Delta}_1)\amalg_{(\mathcal X_0,\unsure{1|0})}\unsure{\mathcal X_1}
\]
\deleted{$\unsure{\mathcal X_0\times\widehat{\Delta}_1}$}}

\deleted{\unsure{Le}} Le ``foncteur recollement'' \inserted{(exact à g.)} (de la théorie générale
\[
\varphi:\mathcal F(\mathcal X_0)\mathbin{-}\mathcal F(\mathcal X_1)
\]
n'est autre que
\[
(4)\qquad\varphi=f_*
\]
